\chapter{Identity, Bootstrapping, and Persistent Process}
\label{ch:identity}

\epigraph{Persistence itself may often be an active maintenance process rather
than a static substance. Identity becomes less like an immutable object and
more like a repeatedly recompiled process maintaining coherence through time.}{}

\section{The Self as Repeated Reconstruction}

A narrative self is not continuously present in the way ordinary intuition
suggests. Continuity is actively reconstructed after sleep, anesthesia, trauma,
distraction, and context switching. The feeling of persistent identity emerges
from this reconstruction process succeeding smoothly enough that the
discontinuities become invisible.

\claimP This is not merely a phenomenological observation but has formal
structure in the trajectory framework. Identity is a trajectory through a state
space that satisfies local compatibility conditions at each step. The compatibility
conditions are the accumulated constraints of prior commitments, relationships,
memories, and bodily continuity. The ``self'' at any moment is the locally
admissible configuration that satisfies all these constraints simultaneously.
When the constraints are temporarily suspended---during sleep, anesthesia,
trauma---the reconstruction process must re-assemble an admissible configuration
from available traces. The resulting configuration is the self after the
interruption.

\section{Computational Bootstrapping as Philosophical Model}

The layered initialization sequence of an operating system provides a structural
model for how complex coherent systems emerge from primitive constraint-stabilization
events.

Firmware initializes primitive hardware assumptions. The bootloader establishes
a minimal execution environment. The kernel initializes memory structures,
scheduling systems, and drivers. Higher-level abstractions---file systems,
user interfaces, application frameworks---only become meaningful because these
lower-level structures have already been stabilized. Each layer depends on the
layers below it having been successfully initialized. Skipping a layer does not
produce a partial system; it produces a system that fails to initialize at all,
because the abstractions of layer $n+1$ are not meaningful without the
constraints of layer $n$.

\claimP The analogy extends to theoretical understanding. Concepts must be
bootstrapped in dependency order to produce genuine comprehension rather than
surface familiarity. A student who memorizes the statement of the second law of
thermodynamics without first developing intuition for entropy, irreversibility,
and heat flow has initialized a terminal abstraction without its boot sequence.
The statement may be reproduced accurately under examination; it cannot be applied
to novel problems because the inferential machinery that generated it is not present.

This is why the monograph's chapter ordering enforces the dependency structure
rather than proceeding from the most striking claims. The Yarncrawler formal
definition in Chapter~\ref{ch:yarncrawler_formal} is not meaningful without the
admissibility and reconstruction vocabulary of Chapters~\ref{ch:manifold}
and~\ref{ch:vocabulary}. The sheaf-theoretic repair machinery of
Chapter~\ref{ch:sheaf} is not meaningful without the Yarncrawler definition.
The cosmological application of Chapter~\ref{ch:cosmological} is not meaningful
without both. Each chapter is a boot layer.

\section{Persistence as Active Maintenance}

Selves, institutions, scientific paradigms, languages, and software ecosystems
persist not by being immutable objects but by continuously reproducing enough
structural continuity to remain operational across interruptions and
transformations.

\claimP A scientific paradigm persists not because its core claims are immune
to revision but because each generation of practitioners reinitializes the
paradigm from available evidence and prior commitments, producing a reconstruction
that is close enough to prior reconstructions that the paradigm remains recognizable.
When the reconstructions begin diverging significantly---when the boot sequence
produces different results in different laboratories---the paradigm is undergoing
revision even if no single paper announces it.

\claimP A language persists not because its words have fixed meanings but because
each generation of speakers reconstructs approximately the same set of
constraint relations among words from available input. The language drifts when
the reconstruction process begins producing systematically different constraint
relations. The drift is invisible at any given moment---each speaker's reconstruction
is locally admissible---and only visible retrospectively when the accumulated
difference becomes too large to ignore.

\section{Semantic, Episodic, and Procedural Memory as Boot Layers}

The computational memory literature distinguishes three types of long-term memory
in agentic systems, a taxonomy borrowed from cognitive psychology:

\begin{itemize}
  \item \textbf{Semantic memory}: generalized factual knowledge, stable across
    contexts. In the trajectory framework, semantic memory encodes the stable
    regions of the admissibility manifold---the constraint structures that do
    not change with individual interaction histories.
  \item \textbf{Episodic memory}: records of specific past interactions and
    their outcomes. This is structural residue in the formal sense of
    Definition~\ref{def:residue}: the accumulated traces of prior constraint
    configurations that continue influencing future trajectory selection.
  \item \textbf{Procedural memory}: the instructions and operational rules that
    govern how the system executes its tasks---the system prompt, the tool-use
    guidelines, the triage criteria.
\end{itemize}

\claimP This trichotomy maps directly onto the boot sequence architecture of
this chapter. Semantic memory is the lowest boot layer: the foundational
constraints that must be initialized before higher-level operations are
meaningful. Episodic memory is the middle layer: the historical record that
personalizes general constraints to a specific interaction context. Procedural
memory is the highest layer: the operational instructions that are only
meaningful given the semantic and episodic layers beneath them.

The crucial observation---which recent agentic systems are beginning to implement
formally---is that procedural memory must be dynamically updatable. A system with
fixed procedural memory (hardcoded system prompts, static triage rules) cannot
adapt to changing constraint conditions. It can only navigate the admissibility
space defined at initialization. As the operational environment changes---as
users provide feedback, as workflow requirements evolve, as edge cases accumulate
in the episodic record---the procedural layer must be revised to maintain
corrigibility.

\subsection{Procedural Memory as Dynamic Admissibility Revision}

In Spherepop terms, updating procedural memory is a Bind operation at the
level of the system's own operational rules: a new dependency is established
between observed feedback (a committed event in the episodic record) and the
future admissibility conditions under which the system will operate. The update
is irreversible in the sense that the feedback has been incorporated; but unlike
Pop or Collapse, the procedural update is designed to be further revisable by
subsequent feedback.

\claimH This is exactly the RC-CLIO repair cycle (Appendix~\ref{app:clio})
applied to the operational layer: seam detection (feedback identifies a mismatch
between current behavior and desired behavior), cocycle identification (the
specific instruction that is producing the mismatch), repair morphism
construction (an LLM generates a revised instruction that would have produced
the correct behavior), and residue annotation (the revision is stored with its
provenance so future revisions can build on it rather than overwriting it blindly).

\subsection{The Provenance Problem in Procedural Updates}

The critical risk in dynamic procedural memory is the same as the proxy
stabilization risk identified in Chapter~\ref{ch:ellul}: the updated
instructions can drift from the underlying operational reality if the feedback
loop is slower than the update rate, or if updates optimize for local feedback
approval rather than global behavioral coherence.

\claimP An LLM that updates system prompts based on user feedback is performing
exactly the proxy drift operation: it is modifying $r$ (the procedural memory)
in response to signals that may not accurately track the target system $s$
(the full space of admissible behaviors). A single piece of negative feedback
may cause an overly broad revision that breaks previously correct behavior.
A pattern of positive feedback may reinforce a locally pleasing but globally
inadmissible trajectory.

The provenance requirement therefore applies with full force to procedural
memory systems. Each update to a system prompt should be stored with its
complete dependency record: which feedback event triggered it, what the
prior instruction was, what alternatives were considered, and what the
predicted behavioral change is. Without this record, a series of individually
reasonable updates can accumulate into a procedural state that no one
designed and no one can fully reconstruct---the institutional identity
proxy problem of Section~\ref{sec:ai_labor} reproduced at the
agent level.

\begin{constraintboundary}
\textbf{Status:~\textbf{[P/H]}} The mapping of semantic/episodic/procedural
memory onto the boot sequence architecture is phenomenological. The RC-CLIO
correspondence is a heuristic engineering interpretation, not a formal
derivation.

\textbf{Design implication:} Procedural memory systems should implement
provenance annotation as a first-class feature, not an afterthought. The
update history of a system prompt is as important as its current state,
for the same reasons that a version control system's commit history is
more informative than the current file contents alone.

\textbf{Falsification:} A procedural memory system that achieves stable,
coherent long-term behavioral adaptation without provenance annotation---
in which updates reliably improve global behavior despite operating only
on compressed local feedback signals---would weaken the provenance
requirement claim.
\end{constraintboundary}

When the maintenance process for identity begins operating on compressed summaries
rather than reconstructed structure, the proxy stabilization failure mode
reappears at the ontological level. The remembered profile, institutional brand,
or theoretical paradigm begins functioning as a gravitational attractor shaping
future interpretation rather than a faithful record enabling genuine reconstruction.

\claimP An institution whose identity has become a proxy---whose members navigate
the institution's self-description rather than its actual operational structure---
has undergone proxy stabilization at the level of collective self-model. The
stated mission, the founding documents, the official history become the proxy
$r$; the actual constraint structure and operational reality are the target $s$;
and optimization pressure (for internal coherence, for external legibility, for
fundraising) drives $r$ and $s$ apart at a rate proportional to the pressure and
inversely proportional to the feedback from actual operations.

The formal structure is identical to the proxy drift of Chapter~\ref{ch:ellul}.
The domain is different: not a platform metric but a self-model. The mechanism
is the same.

\begin{constraintboundary}
\textbf{Established:} The formal connection between identity and trajectory
maintenance under constraint. The bootstrapping model as an account of how
higher-level cognitive and institutional structures depend on lower-level
initialization.

\textbf{Phenomenological~[P]:} The applications to scientific paradigms,
language drift, and institutional identity. These are interpretive frameworks
supported by the formal structure but not derived from it.

\textbf{Not claimed:} That biological identity is literally a computational
process. The computational bootstrapping analogy is structural, not ontological.
The analogy licenses operational insight transfer insofar as the relevant
constraint structures are analogous; it does not imply that minds are computers.
\end{constraintboundary}
